% =====================================================================
%  FICHE 11 — THEME 1 — Corrigé FACILE — Le nombre d'oxydation
% =====================================================================
\documentclass[11pt,a4paper]{article}
\usepackage[utf8]{inputenc}
\usepackage[T1]{fontenc}
\usepackage{lmodern}
\usepackage[french]{babel}
\usepackage{amsmath,amssymb}
\usepackage[version=4]{mhchem}
\usepackage{siunitx}
\sisetup{output-decimal-marker={,},per-mode=symbol}
\usepackage{xcolor}
\usepackage{array}
\usepackage{booktabs}
\usepackage{enumitem}
\usepackage[most]{tcolorbox}
\usepackage[a4paper,margin=2.1cm,headheight=26pt,headsep=10pt]{geometry}
\usepackage{fancyhdr}
\usepackage[hidelinks]{hyperref}

\definecolor{primary}{RGB}{146,95,160}
\definecolor{primarydark}{RGB}{92,52,104}
\definecolor{corrcol}{RGB}{0,135,90}
\newcommand{\NO}{\ensuremath{\mathrm{N.O.}}}

\newcounter{exo}
\newtcolorbox{corr}[1][]{enhanced,breakable,boxrule=0.9pt,arc=2pt,colback=corrcol!4,
  colframe=corrcol,fonttitle=\bfseries,coltitle=white,left=3mm,right=3mm,top=2.2mm,bottom=2.2mm,
  title={Corrigé~\refstepcounter{exo}\theexo\ifx\relax#1\relax\relax\else\ \textmd{--- \itshape #1}\fi}}

\pagestyle{fancy}\fancyhf{}
\fancyhead[L]{\small\textcolor{primarydark}{\textbf{Fiche 11 · Thème 1} — Corrigé}}
\fancyhead[R]{\small\textcolor{primarydark}{\textsc{Facile}}}
\fancyfoot[C]{\small\thepage}
\renewcommand{\headrulewidth}{0.8pt}
\renewcommand{\headrule}{\hbox to\headwidth{\color{primary}\leaders\hrule height \headrulewidth\hfill}}

\begin{document}
\begin{center}
{\LARGE\bfseries\color{primarydark}Thème 1 — Corrigé Facile}\\[-2pt]
{\color{corrcol}\rule{\textwidth}{1pt}}
\end{center}

\begin{corr}[corps simples et ions monoatomiques]
\begin{itemize}[leftmargin=1.5em,topsep=1pt,itemsep=1.5pt]
  \item \ce{O2} : corps simple (règle R1) $\Rightarrow \NO(\ce{O})=0$.
  \item \ce{Na+} : ion monoatomique de charge $+1$ (R2) $\Rightarrow \NO(\ce{Na})=+\mathrm{I}$.
  \item \ce{Al^{3+}} : ion monoatomique de charge $+3$ (R2) $\Rightarrow \NO(\ce{Al})=+\mathrm{III}$.
  \item \ce{Cl^-} : ion monoatomique de charge $-1$ (R2) $\Rightarrow \NO(\ce{Cl})=-\mathrm{I}$.
  \item \ce{Fe(s)} : corps simple, métal à l'état natif (R1) $\Rightarrow \NO(\ce{Fe})=0$.
  \item \ce{P4} : corps simple (R1) $\Rightarrow \NO(\ce{P})=0$.
\end{itemize}
\textbf{Retenir :} corps simple $\to 0$ ; ion monoatomique $\to$ sa charge. Aucun calcul.
\end{corr}

\begin{corr}[le \NO\ de l'oxygène]
\begin{itemize}[leftmargin=1.5em,topsep=1pt,itemsep=1.5pt]
  \item \ce{CaO} : cas général $\Rightarrow \NO(\ce{O})=-\mathrm{II}$ (on vérifie : $\NO(\ce{Ca})=+\mathrm{II}$, somme $=0$).
  \item \ce{CO2} : cas général $\Rightarrow \NO(\ce{O})=-\mathrm{II}$.
  \item \ce{H2O2} : \textbf{peroxyde} (liaison \ce{O-O}), exception R4 $\Rightarrow \NO(\ce{O})=-\mathrm{I}$
        (vérif. : $2(+1)+2x=0\Rightarrow x=-1$).
  \item \ce{Na2O2} : \textbf{peroxyde} de sodium, exception R4 $\Rightarrow \NO(\ce{O})=-\mathrm{I}$
        (vérif. : $2(+1)+2x=0\Rightarrow x=-1$).
\end{itemize}
\textbf{Piège classique :} reconnaître un peroxyde ($-\mathrm{I}$) et non un oxyde ordinaire ($-\mathrm{II}$).
\end{corr}

\begin{corr}[le \NO\ de l'hydrogène]
\begin{itemize}[leftmargin=1.5em,topsep=1pt,itemsep=1.5pt]
  \item \ce{HCl} : cas général $\Rightarrow \NO(\ce{H})=+\mathrm{I}$.
  \item \ce{NH3} : cas général $\Rightarrow \NO(\ce{H})=+\mathrm{I}$.
  \item \ce{CaH2} : \textbf{hydrure métallique} (exception R3) $\Rightarrow \NO(\ce{H})=-\mathrm{I}$.
  \item \ce{H2} : corps simple (R1) $\Rightarrow \NO(\ce{H})=0$.
\end{itemize}
L'\textbf{hydrure} est \ce{CaH2} : le calcium (métal) est moins électronégatif que l'hydrogène, qui
joue alors le rôle d'anion hydrure \ce{H^-}. On vérifie : $\NO(\ce{Ca})=+\mathrm{II}$ et
$(+\mathrm{II})+2(-\mathrm{I})=0$.
\end{corr}

\begin{corr}[un seul atome inconnu]
On écrit à chaque fois $\sum\NO=0$ (molécules neutres), avec $\NO(\ce{H})=+\mathrm{I}$ et $\NO(\ce{O})=-\mathrm{II}$.
\begin{enumerate}[label=(\alph*),leftmargin=2em,topsep=1pt,itemsep=1.5pt]
  \item \ce{SO2} : $x+2(-2)=0 \Rightarrow x=+4$, donc $\NO(\ce{S})=+\mathrm{IV}$.
  \item \ce{CH4} : $x+4(+1)=0 \Rightarrow x=-4$, donc $\NO(\ce{C})=-\mathrm{IV}$.
  \item \ce{NH3} : $x+3(+1)=0 \Rightarrow x=-3$, donc $\NO(\ce{N})=-\mathrm{III}$.
\end{enumerate}
\end{corr}

\begin{corr}[vérifier une attribution]
Avec $\NO(\ce{S})=+\mathrm{IV}$ et $\NO(\ce{O})=-\mathrm{II}$, et 1~atome de S pour 2~atomes de O :
\[
\sum\NO = (+4) + 2\times(-2) = +4 - 4 = 0.
\]
La somme vaut bien $0$, ce qui est \textbf{cohérent} avec le caractère neutre de \ce{SO2}.
L'attribution proposée est donc correcte. \emph{Toujours utiliser cette égalité comme contrôle final.}
\end{corr}

\end{document}
